\lecture{17}{Ньютоновский поток и дивергенция Брэгмана}

\sourcevideo
    {https://www.youtube.com/playlist?list=PLOzRYVm0a65fhKDS8cYIC7YCuVGJ4FOJx}
    {Week 5: Lecture 17: Bregman Divergence}

Обычный градиентный поток использует только направление
\(-\nabla f(x)\) и потому чувствителен к различию кривизны по
координатам. Ньютоновский поток компенсирует эту анизотропию обратным
Гессианом, а дивергенция Брэгмана заменяет евклидово расстояние мерой
расхождения, согласованной с геометрией задачи.

\section{Строгая выпуклость и скорость градиентного потока}

Пусть \(f\in C^2(\RR^n)\) и
\[
    \nabla^2f(x)\succ0
\]
во всей выпуклой области.

\begin{theorem}
Если Гессиан положительно определён в каждой точке выпуклой области,
то функция \(f\) строго выпукла на этой области.
\end{theorem}

Обратное утверждение неверно. Функция
\[
    f(x)=x^4
\]
строго выпукла, но
\[
    f''(0)=0.
\]
Строгая выпуклость гарантирует единственность минимума, если он
существует, но не гарантирует его существование. Например,
\(f(x)=e^x\) строго выпукла на \(\RR\), однако не достигает своего
инфимума.

Сильная выпуклость требует единой оценки
\[
    \nabla^2f(x)\succeq\mu I,
    \qquad
    \mu>0.
\]
При одной строгой выпуклости такая константа может отсутствовать.
Для \(f(x)=x^4\) градиентный поток
\[
    \dot x=-4x^3
\]
имеет при \(x(0)=x_0\ne0\) решение
\[
    x(t)
    =
    \frac{x_0}{\sqrt{1+8x_0^2t}},
\]
поэтому
\[
    |x(t)|=O(t^{-1/2}).
\]
Итак, строгая выпуклость обеспечивает единственность решения, но не
равномерную экспоненциальную скорость обычного градиентного потока.

\section{Ньютоновский поток}

Предположим, что
\[
    \nabla^2f(x)\succ0
\]
для всех рассматриваемых \(x\). Тогда можно определить динамику
\[
    \dot x
    =
    -\bigl(\nabla^2f(x)\bigr)^{-1}\nabla f(x).
\]
Она называется непрерывным методом Ньютона, или ньютоновским
потоком. В дискретном времени той же идее соответствует шаг
\[
    x^{k+1}
    =
    x^k-
    \bigl(\nabla^2f(x^k)\bigr)^{-1}\nabla f(x^k).
\]

Рассмотрим функцию Ляпунова
\[
    V(x)=\frac12\lVert\nabla f(x)\rVert^2.
\]
Если \(x^\star\) — единственная стационарная точка, то \(V\)
положительно определена относительно \(x^\star\). Вдоль
ньютоновского потока
\[
\begin{aligned}
    \dot V
    &=
    \nabla f(x)^{\mathsf T}\nabla^2f(x)\dot x
    \\
    &=
    -\nabla f(x)^{\mathsf T}\nabla^2f(x)
    \bigl(\nabla^2f(x)\bigr)^{-1}\nabla f(x)
    \\
    &=
    -\lVert\nabla f(x)\rVert^2
    =-2V.
\end{aligned}
\]
Следовательно,
\[
    V(x(t))
    =
    e^{-2(t-t_0)}V(x(t_0))
\]
и
\[
    \lVert\nabla f(x(t))\rVert
    =
    e^{-(t-t_0)}
    \lVert\nabla f(x(t_0))\rVert.
\]

Тот же результат получается непосредственно. Для
\[
    g(t)=\nabla f(x(t))
\]
имеем
\[
    \dot g
    =
    \nabla^2f(x)\dot x
    =-g,
\]
то есть нелинейная динамика состояния превращается в линейную
динамику градиента.

Экспоненциальное убывание градиента не автоматически даёт такую же
оценку расстояния до минимума. Для перехода нужна связь вида
\[
    \lVert\nabla f(x)\rVert
    \ge
    \mu\lVert x-x^\star\rVert.
\]
Она выполняется при \(\mu\)-сильной выпуклости, но может отсутствовать
при одной строгой выпуклости. Кроме того, для ньютоновского потока
нужно отдельно обеспечить существование траектории и невырожденность
Гессиана вдоль неё.

Если \(f\) только выпукла и
\[
    \nabla^2f(x)\succeq0,
\]
обратный Гессиан может не существовать. Для обычного градиентного
потока
\[
    \dot x=-\nabla f(x)
\]
та же функция даёт
\[
    \dot V
    =
    -\nabla f(x)^{\mathsf T}\nabla^2f(x)\nabla f(x)
    \le0.
\]
Отсюда следует невозрастание нормы градиента, но не скорость
сходимости и не обязательно сходимость к заранее выбранному
минимизатору. При нескольких решениях \(V\) положительно определена
относительно всего множества стационарных точек.

\section{Компенсация плохой обусловленности}

Рассмотрим квадратичную функцию
\[
    f(x,y)
    =
    \frac12x^2+
    \frac{10^{-4}}2y^2.
\]
Её градиент и Гессиан равны
\[
    \nabla f(x,y)
    =
    \begin{pmatrix}
        x\\
        10^{-4}y
    \end{pmatrix},
    \qquad
    H
    =
    \begin{pmatrix}
        1&0\\
        0&10^{-4}
    \end{pmatrix}.
\]
Число обусловленности составляет
\[
    \kappa(H)=10^4.
\]
Обычный градиентный поток имеет вид
\[
    \dot x=-x,
    \qquad
    \dot y=-10^{-4}y,
\]
поэтому
\[
    x(t)=e^{-t}x(0),
    \qquad
    y(t)=e^{-10^{-4}t}y(0).
\]
Медленная координата определяет общую скорость.

Поскольку
\[
    H^{-1}
    =
    \begin{pmatrix}
        1&0\\
        0&10^4
    \end{pmatrix},
\]
ньютоновский поток становится
\[
    \begin{pmatrix}
        \dot x\\
        \dot y
    \end{pmatrix}
    =
    -H^{-1}\nabla f(x,y)
    =
    -
    \begin{pmatrix}
        x\\
        y
    \end{pmatrix}.
\]
Обе координаты затухают как \(e^{-t}\). Обратный Гессиан выравнивает
различную кривизну направлений и преобразует вытянутую квадратичную
геометрию в изотропную.

Нормированный поток
\[
    \dot x
    =
    -\frac{\nabla f(x)}
    {\lVert\nabla f(x)\rVert+\varepsilon},
    \qquad
    \varepsilon>0,
\]
изменяет только длину градиентного вектора и сохраняет его
направление. Обратный Гессиан, напротив, масштабирует координаты
по-разному и в общем случае меняет также направление. Поэтому
скалярная нормировка не устраняет анизотропию так, как это делает
матричное преобразование второго порядка.

Основные евклидовы функции Ляпунова в оптимизации равны
\[
    f(x)-f^\star,
    \qquad
    \frac12\lVert x-x^\star\rVert^2,
    \qquad
    \frac12\lVert\nabla f(x)\rVert^2.
\]
Первая обращается в нуль на множестве глобальных минимумов, вторая
— только в выбранной точке \(x^\star\), третья — на множестве
стационарных точек. Для невыпуклой функции стационарность не
равносильна глобальной оптимальности.

\section{Дивергенция Брэгмана}

Евклидова норма не всегда соответствует структуре задачи. Пусть
\(h\colon\mathcal D\to\RR\) дифференцируема на выпуклом множестве
\(\mathcal D\).

\begin{definition}[Дивергенция Брэгмана]
Дивергенцией Брэгмана, порождённой функцией \(h\), называется
\[
    D_h(p,q)
    =
    h(p)-h(q)
    -\langle\nabla h(q),p-q\rangle.
\]
\end{definition}

Величина
\[
    h(q)+\langle\nabla h(q),p-q\rangle
\]
является аффинной касательной моделью \(h\) в точке \(q\),
вычисленной в \(p\). Поэтому \(D_h(p,q)\) измеряет ошибку этой
линейной аппроксимации.

Из выпуклости первого порядка следует
\[
    D_h(p,q)\ge0.
\]
Если \(h\) строго выпукла, то
\[
    D_h(p,q)=0
    \quad\Longleftrightarrow\quad
    p=q.
\]
При одной выпуклости нулевое значение возможно и для различных точек,
например на аффинном участке функции.

Если \(h\) \(\mu\)-сильно выпукла, то
\[
    D_h(p,q)
    \ge
    \frac\mu2\lVert p-q\rVert^2.
\]
Если она дополнительно \(L\)-гладка, то
\[
    D_h(p,q)
    \le
    \frac L2\lVert p-q\rVert^2.
\]
Таким образом,
\[
    \frac\mu2\lVert p-q\rVert^2
    \le
    D_h(p,q)
    \le
    \frac L2\lVert p-q\rVert^2.
\]
Сильная выпуклость даёт нижнюю оценку, а не точное равенство.
Равенство возникает лишь для специальных квадратичных функций.

При
\[
    h(x)=\frac12\lVert x\rVert^2
\]
получаем
\[
    D_h(p,q)=\frac12\lVert p-q\rVert^2.
\]
Для
\[
    h(x)=\frac12x^{\mathsf T}Px,
    \qquad
    P\succ0,
\]
имеем
\[
    D_h(p,q)
    =
    \frac12(p-q)^{\mathsf T}P(p-q).
\]
Матрица \(P\) задаёт взвешенную евклидову геометрию.

В общем случае дивергенция Брэгмана не является метрикой:
\[
    D_h(p,q)\ne D_h(q,p)
\]
и не обязана удовлетворять неравенству треугольника. Полезным
заменителем евклидовых тождеств служит трёхточечная формула
\[
\begin{aligned}
    D_h(p,q)+D_h(q,r)-D_h(p,r)
    =
    \langle\nabla h(r)-\nabla h(q),p-q\rangle.
\end{aligned}
\]
Она широко используется в анализе зеркальных методов.

\section{Брэгмановская геометрия в алгоритмах}

Если \(x^\star\) — оптимальное решение, то естественными
кандидатами на функцию Ляпунова являются
\[
    D_h(x^\star,x)
    \qquad\text{и}\qquad
    D_h(x,x^\star).
\]
Порядок аргументов существенен из-за несимметричности. При строгой
выпуклости \(h\) обе величины неотрицательны и обращаются в нуль
только при \(x=x^\star\); при сильной выпуклости они дополнительно
контролируют квадрат евклидова расстояния. Конкретная производная
вдоль траектории зависит от выбранного потока, поэтому функцию \(h\)
и динамику нужно проектировать совместно.

Евклидов градиентный шаг можно записать как
\[
    x^{k+1}
    =
    \operatorname*{argmin}_{x}
    \left\{
        \langle\nabla f(x^k),x-x^k\rangle
        +
        \frac1{2\eta}\lVert x-x^k\rVert^2
    \right\}.
\]
Замена квадрата нормы дивергенцией Брэгмана даёт зеркальный шаг
\[
    x^{k+1}
    =
    \operatorname*{argmin}_{x}
    \left\{
        \langle\nabla f(x^k),x-x^k\rangle
        +
        \frac1\eta D_h(x,x^k)
    \right\}.
\]
Функция \(h\) определяет геометрию регуляризующего члена и может быть
согласована с допустимым множеством и масштабом координат. Зеркальный
спуск подробно в лекции не выводится; эта запись поясняет практическую
роль дивергенции.

Квадратичная \(h\) приводит к взвешенной евклидовой геометрии, а
неквадратичная позволяет строить существенно неевклидовы шаги.
Вариационный взгляд на ускоренные методы использует подобные
дивергенции для построения непрерывных динамик и функций энергии. 

\section{Вычислительная стоимость и границы метода}

Обычный градиентный поток требует вычисления \(\nabla f(x)\).
Ньютоновский поток дополнительно требует Гессиан и решение системы
\[
    \nabla^2f(x)d=-\nabla f(x).
\]
Явно строить обратную матрицу не следует. Для плотного Гессиана
размерности \(n\) факторизация обычно требует порядка
\[
    O(n^3)
\]
операций и
\[
    O(n^2)
\]
памяти. Поэтому методы второго порядка могут использовать меньше
итераций, но каждая итерация существенно дороже.

Брэгмановские методы обычно не требуют Гессиана целевой функции, но
требуют решения вспомогательной задачи, определяемой \(h\). Их
преимущество проявляется, когда евклидова геометрия плохо отражает
структуру пространства или ограничений.

Ньютоновский поток особенно полезен при обратимом Гессиане, сильной
анизотропии и доступном решении линейных систем. Нормированный
градиент управляет общей длиной шага, но не выравнивает координатную
кривизну. Дивергенция Брэгмана позволяет согласовать функцию
Ляпунова и регуляризующий член с неевклидовой геометрией задачи.
